\chapter*{\centerline{Abstract}}
\markboth{\MakeUppercase{Abstract}}{}
\iftoggle{fulltoc}{
  \addcontentsline{toc}{chapter}{Abstract}
}{}

Biomimetic robots apply biological knowledge to robotics and enable controlled experiments in biomechanics and neuroscience. This dissertation addresses four needs for a human-proportioned robot with braided pneumatic actuators (BPAs): force prediction across resting lengths, corrected joint-torque predictions, accessible sensory-afferent evidence, and tools for transferring and testing neural controllers between biomechanical models. Festo BPA characterization across resting lengths, contractions, and pressures produced a resting-length-dependent maximum-force model and a three-coefficient nondimensional force surface. Joint-torque design factors account for length offsets, effective series compliance, and wrapping-induced loss of usable actuator length. Redesigned paired BPA flexor and extensor routes met or exceeded selected human-muscle reference torques at the measured knee angles. The Sensory Afferent Database (SADb) connects experimental findings and model assumptions through searchable notes and citation navigation. Neural-model tools, graphical blocks, and simulation interfaces enable controlled experiments in MuJoCo and Simulink/Simscape. Evaluations establish selected circuit responses, implementation agreement, and virtual-harness-assisted balance and walking behaviors. Together, these contributions support experimentally informed robot design and neuromechanical experiments, with potential prosthetic and orthotic applications.
